\begin{figure*}[t] % fig:03-framework-viz-synthesis
    \centering
    \includegraphics[width=\textwidth]{./figures/03-framework-viz-synthesis.pdf}
    \caption{
        The Synthesis pipeline to generate our image priors.
        It includes 3 steps: (1) We initialize hierachical noises $x_{\text{hn}}$; (2) We apply non-linear transforms to create synthetic images $x_{\text{nt}}$; and (3) We construct the final images $\mathcal{X}_{\mathcal{P}}$ by applying semantic cutmixing with a semantic mask $\hat{m}$.
        The details of each step is describe in text.
        Compared to noise images sampled from random Gaussian distribution $\mathcal{X}_{\mathcal{N}}$ and random uniform distribution $\mathcal{X}_{\mathcal{U}}$, our image priors $\mathcal{X}_{\mathcal{P}}$ have much more diverse and distinct patterns.
    }
    \label{fig:03-framework-viz-synthesis}
\end{figure*}

\textbf{Problem statement.}
\begin{foldable} % Problem statement
    Given solely a \textit{black-box} pre-trained teacher network $T$ that returns only top-1 hard labels and \textit{no training data}, our goal is to train a student network $S$ to approximate $T$.
\end{foldable}

\begin{foldable} % Naive solution
    Lacking a direct solution, we may tweak the standard KD framework~\cite{hinton2015distilling} to our setting.
    Without data, we resort to using noise images, \eg, $x\sim\mathcal{X}_{\mathcal{U}}$, where $\mathcal{X}_{\mathcal{U}}=\mathcal{U}\left[0,1\right]^{C\times H\times W}$ and $C\times H\times W$ is the image dimension.
    Moreover, we can only use a single hard label $\hat{y}_{T-\text{Hard}}=T(x)\in\left\{ 1,\ldots,K\right\}$ to match the student's probabilistic outputs $\hat{y}_{S-\text{Soft}}=S\left(x\right)$ as a naive solution:
    \begin{equation} \label{eq:loss-naive-kd}
        \mathcal{L}_{\text{Naive-KD}}=\mathbb{E}_{x\sim\mathcal{X}_{\mathcal{U}}}\left[\mathcal{L}_{\text{CE}}\left(\hat{y}_{S-\text{Soft}},\hat{y}_{T-\text{Hard}}\right)\right],
    \end{equation}
    where $\mathcal{L}_{\text{CE}}$ is the cross-entropy loss.
    While this approach can handle simple cases, we will show that it highly depends on the knowledge captured in the synthetic images and fails to use inter-class knowledge --- the pivotal advantage in KD.
    To circumvent this challenge, previous methods~\cite{wang2021zero, zhang2022ideal, yuan2024data} have had different approaches, but none directly alleviates the diversity concern.
\end{foldable}

\textbf{Proposal.}
\begin{foldable} % Proposal
    To improve image diversity in BBDFKD, we propose our novel framework \textbf{\methodname} of three components: Synthesis, Contrastion, and Distillation.
    In the \textit{Synthesis} phase, we design the crafting of \textit{image priors} $\mathcal{X}_\mathcal{P}$ --- synthetic images that are compounded with hierarchy-level diversity, nonlinearity, and meaningful semantics.
    Our goal is to capture the complex and diverse patterns in natural scenes and objects.
    Then, in the \textit{Contrastion} phase, we further optimize images $\hat{x}\sim\mathcal{X}_\mathcal{P}$ for diversity with contrastive learning, to construct a synthetic dataset $\mathcal{D}$.
    Finally, in the \textit{Distillation} phase, we use $\mathcal{D}$ to distill the student $S$ from the teacher $T$.
\end{foldable}


\subsection{Synthesis: crafting diverse image priors} \label{ssec:03-framework-synthesis}
\begin{foldable} % Motivation --> overview of 3 components
    From a visual recognition point-of-view, there exists \textit{`strong local structures'} between neighbor pixels that compose \textit{global patterns}~\cite{lecun2002gradient} in natural scenes and objects.
    Centered around this principle, we create image priors $\mathcal{X}_\mathcal{P}$ with our Synthesis pipeline of three modules: \textit{hierarchical noise}, \textit{nonlinear transformation}, and \textit{semantic cutmixing}.
    We visualize the pipeline in Fig.~\ref{fig:03-framework-viz-synthesis} and analyze their image representations in the experiments (Sec.~\ref{sssec:04-experiments-ablation-embeddings-generation}).
\end{foldable}

\subsubsection{Hierarchical noise} \label{sssec:03-framework-synthesis-hiernoise}
\begin{foldable} % Hierchical structure in nature
    Akin to how a car is characterized by its windows, mirrors, wheels, an animal is easily recognizable via its eyes, ears, limbs.
    Hierarchy is an universal property transcending human perception.
    Based on this observation, we design a sampling technique that captures image patterns at both the local and global scales to mimic the hierarchical structures often seen in natural scenes and objects.
\end{foldable}

\begin{foldable} % Sampling multi-scale noise images
    For a pure pixel-wise noise image $x\sim\mathcal{X_{\mathcal{U}}}$, the patterns are majorly local and independent.
    Whereas, a monochromatic image (\ie, single-color) has dominant global correlation.
    Along this spectrum, there must exist a trade-off between the local-versus-global property.
    Inspiredly, we design a sampler composing noise images at multiple scales.
    Let us sample a collection of noise images with increasing dimensions until they sufficiently envelope $H\times W$:
    \begin{equation} \label{eq:synthesis-hiernoise-multiscale}
        \left\{ x_{d}|x_{d}\sim\mathcal{U}\left[0,1\right]^{C\times2^{d}\times2^{d}}\right\} _{d=0}^{d_{\text{max}}},
    \end{equation}
    where $d_{\text{max}}=\text{ceil}\left(\log_{2}\max\left(H,W\right)\right)$ is the sufficient scale.
\end{foldable}

\begin{foldable} % Combining multi-scale noise images
    Next, we aim to compose this collection of local-to-global noise to a single image.
    Thus, we proceed to sample coefficients $\displaystyle \left\{ \alpha_{d}|\alpha_{d}\sim\mathcal{N}\left(0,1\right) \right\}_{d=0}^{d_{\text{max}}}$ and put them through softmax to obtain each $\displaystyle \bar{\alpha}_{d}=\frac{\exp\left(\alpha_{d}\right)}{\sum_{d^\prime=0}^{d_\text{max}}\exp\left(\alpha_{d^\prime}\right)}$.
    This ensures $\sum_{d=0}^{d_{\text{max}}}\bar{\alpha}_{d}=1$ and $\bar{\alpha}_{d}~\in~[0, 1]$ for a valid combination, which is our \textit{hierarchical noise image}:
    \begin{equation} \label{eq:synthesis-hiernoise-combine}
        x_{\text{hn}}=\sum_{d=0}^{d_{\text{max}}}\bar{\alpha}_{d}\cdot f_{\text{upscale}}\left(x_{d};2^{d_\text{max}}\right),
    \end{equation}
    where $f_{\text{upscale}}\left(\cdot;2^{d_{\text{max}}}\right)$ is an upscaling transformation so that noise images $\left\{ x_{d}\right\}$  share equal dimensions to allow additivity.
\end{foldable}

\subsubsection{Nonlinear transformation} \label{sssec:03-framework-synthesis-nonlintrans}
\begin{foldable} % Nonlinear transformation
    To enrich nonlinear patterns of synthetic images, we apply random rotation $f_{\text{rot}}$ of $\left[-45^{\circ},45^{\circ}\right]$, random elastic transform $f_{\text{elas}}$, and random cropping $f_{\text{crop}}$ to the specified size $H\times W$.
    These operations form our \textit{nonlinear-transformed noise image}:
    \begin{equation} \label{eq:synthesis-nonlintrans}
        x_{\text{nt}}=f_{\text{crop}}^{H\times W} \circ f_{\text{elas}} \circ f_{\text{rot}}^{\left[-45^{\circ},45^{\circ}\right]}\left(x_{\text{hn}}\right).
    \end{equation}
\end{foldable}

\subsubsection{Semantic cutmixing} \label{sssec:03-framework-synthesis-semcutmix}
\begin{foldable} % Motivation
    Finally, the semantics of natural objects usually consist of distinct, cohesive shapes as the 'content' and are distinguished with the background.
    Inspired by CutMix~\cite{yun2019cutmix}, we propose an improved masking technique to cut-mix any two images together, creating diverse and semantically cohesive shapes.
\end{foldable}

\begin{foldable} % Design: diverging filter, mask, and cutmix
    We capture the semantics within a mask $m$, initialized to $m_{0}\sim\mathcal{U}\left[0,1\right]^{1\times H\times W}$.
    Regarding where $m<\beta$ as negative regions and $m\ge\beta$ as positive regions, the key to creating structurally cohesive shapes is to connect same-type regions, but contrast opposite-type ones.
    While the former can be easily executed with a blurring filter $f_{\text{blur}}$, we design the diverging filter $f_{\text{divg}}$ for images within the $\left[0,1\right]$ dynamic range as:
    \begin{equation} \label{eq:synthesis-diverging-filter}
        f_{\text{divg}}\left(m;0\le\beta\le1\right)=\begin{cases}
            \frac{m^{2}}{\beta}            & ,x\ge\beta \\
            \frac{m^{2}-2x+\beta}{\beta-1} & ,x<\beta
        \end{cases}.
    \end{equation}
    We provide detailed formulation of $f_{\text{divg}}$ in Appendix Sec.~\ref{sssec:app01-teacherstudents-optimization-divgfilter}. In hindsight, $f_{\text{divg}}$ is composed of two continuously differentiable half-parabolas: the left-side convex one on $[0,\beta)$, the right-side concave on $\left[\beta,1\right]$, connected at the inflection point $\left(\beta,\beta\right)$.
    To balance the divergence of the negative and positive regions, we select $\beta=0.5$.
    We iteratively (10 times) pass the mask $m_0$ through the two filters as $m_{i+1}=f_{\text{blur}} \circ f_{\text{divg}}\left(m_{i};\beta\right)$ to obtain the final semantic mask $\hat{m}$.
    While our filters take care of creating cohesive shapes, the randomness in $m_{0}$ yields diverse $\hat{m}$.
    Based on $\hat{m}$, we pairwisely cut-mix $\left\{ x_{\text{nt}}\right\}$ against themselves and random monochromatic images $\mathcal{X}_{\text{mc}}=\left\{f_{\text{upscale}}^{H\times W}\left(x_{\text{mc}}\right)|x_{\text{mc}}\sim\mathcal{U}\left[0,1\right]^{C\times1\times1}\right\}$ to form our final \textit{image~priors} $\mathcal{X}_\mathcal{P} = \{\hat{x}\}$:
    \begin{equation} \label{eq:synthesis-semcutmix}
        \hat{x}=\text{CutMix}\left(x_{i},x_{j};\hat{m}|x_{i}\ne x_{j}; x_{i},x_{j}\in\left\{ x_{\text{nt}}\right\} \cup\mathcal{X}_{\text{\text{mc}}}\right).
    \end{equation}
\end{foldable}

\begin{figure}[ht] % fig:03-framework-teaser
    \centering
    \includegraphics[width=\columnwidth]{./figures/03-framework-teaser.pdf}
    \caption{
        Illustration of our method {\methodname} of three phases.
        \textbf{(a) Synthesis:} We sample our image priors $\hat{x}\sim\mathcal{X}_{\mathcal{P}}$
        \textbf{(b) Contrastion:} We construct a dataset $\mathcal{D}=\left\{\hat{x}|\hat{x}\sim\mathcal{X}_{\mathcal{P}}\right\}$ to train an initial student $S_0$ as a feature extractor. Then, we enforce an instance-discriminator $R$ to distinguish embeddings of images $x\sim\mathcal{D}_{g}$ to be similar to its positive view $x^{+}$, and dissimilar to its negative view $x^{-}\ne x$.
        Thus, $\mathcal{D}$ is also optimized to be more diverse.
        \textbf{(c) Distillation:} With diverse $\mathcal{D}$, we distill the final student $S$ by matching with both hard labels from teacher $T$ via $\mathcal{L}_{\text{Hard-KD}}$ and soft labels from $S_0$ via $\mathcal{L}_{\text{Soft-KD}}$.
        More detailed formulations can be found in the texts.
    }
    \label{fig:03-framework-teaser}
\end{figure}


\subsection{Contrastion: improving diversity via contrastive learning} \label{ssec:03-framework-contrastion}
\begin{foldable} % Initial student S_0
    We further improve the diversity of our image priors through contrastive learning.
    However, contrastive methods usually requires extracting the intermediate features, which is not feasible from a black-box teacher.
    To overcome this problem, we simply train an initial student network $S_{0}$ to serve as a foundation for contrastion.
    To train $S_{0}$, we sample the image priors from our Synthesis pipeline to construct an initial dataset $\mathcal{D}=\left\{ x_{i}|x_{i}\sim\mathcal{X}_{\mathcal{P}}\right\} _{i=1}^{N}$.
    We also attempt to make $\mathcal{D}$ class-balanced via rejection sampling, which is efficient as $\mathcal{X}_\mathcal{P}$ is diverse.
    Their labels can be obtained via querying $\hat{y}_{T-\text{Hard}}=T\left(x\right)$.
    We train $S_{0}$ on $\mathcal{D}$ to the following objective:
    \begin{equation} \label{eq:loss-hard-kd}
        \mathcal{L}_{\text{Hard-KD}}=\mathbb{E}_{x\sim\mathcal{D}}\left[\mathcal{L}_{\text{CE}}\left(\hat{y}_{S_0-\text{Soft}},\hat{y}_{T-\text{Hard}}\right)\right],
    \end{equation}
    where $\hat{y}_{S_0-\text{Soft}} = S_0(x)$ is the soft label of $S_{0}$.
    Equipped with a white-box $S_0$, we can explore its rich representations to improve diversity.
\end{foldable}

\begin{foldable} % Embeddings & cosine similarity
    Following~\cite{fang2021contrastive}, diversity boils down to \textit{``how distinguishable are the samples from the dataset''}.
    Interestingly, it can be explicitly formulated as an optimizable objective to mitigate the mode collapse and boost KD performance.
    For this goal, we utilize the student backbone $S_{0}^{\mathcal{H}}$ to extract image embeddings and introduce an \textit{instance-discriminator} $R$ to treat each embedding as a distinct class.
    We argue that even if $S_0$ might not have comparable performance to $T$, its backbone $S_{0}^{\mathcal{H}}$ can still capture the early representations, so $R$ can distinguish them effectively.
    For any two images $(x_i, x_j)$, we quantify their cosine similarity of embeddings as:
    \begin{equation}
        \text{Sim}\left(x_i, x_j\right)=
            \frac{ \left\langle R \circ S_{0}^{\mathcal{H}}\left(x_i\right), R \circ S_{0}^{\mathcal{H}}\left(x_j\right)\right\rangle }
            { \left\Vert R \circ S_{0}^{\mathcal{H}}\left(x_i\right) \right\Vert \cdot\left\Vert R \circ S_{0}^{\mathcal{H}}\left(x_j\right) \right\Vert }.
    \end{equation}
\end{foldable}

\begin{foldable} % Contrastive loss
    To enforce contrastive learning, from a synthetic image $x\sim\mathcal{D}$, we create positive views $x^{+}$ with random augmentation, and sample negative views $x^{-}\sim\mathcal{D}$ different from $x$.
    The contrastive loss is formulated as:
    \begin{equation} \label{eq:loss-contrastive}
        \mathcal{L}_{\text{CT}} = \mathbb{E}_{x \sim \mathcal{D}} \left[ -\log \frac{ \exp\left( \text{Sim}(x, x^{+}) \right) }{ \sum_{x^{-} \ne x} \exp\left( \text{Sim}(x, x^{-})  \right) } \right].
    \end{equation}
    We optimize ate $\mathcal{D}$ w.r.t the objective $\mathcal{L}_{\text{CT}}$ to update it with improved diversity, and use it as our \textit{distillation dataset} .
    % Noticeably, previous techniques~\cite{zhang2022ideal, yuan2024data} employ in-class similarity and class-balance objectives, however, they could cause overfitting, reducing diversity and KD performance~\cite{yuan2024data}.
    % We also experimented with these objectives in Appendix Sec.~\ref{ssec:app02-experiments-multiobjective} but did not have any fruitful results, confirming their ineffectiveness.
\end{foldable}



\subsection{Distillation: Training the student network} \label{ssec:03-framework-distillation}
\begin{foldable}
    After retrieving a distillation dataset $\mathcal{D}$ from the Contrastion step, we train the student $S$.
    For this goal, existing BBDFKD methods~\cite{zhang2022ideal, yuan2024data} only use the Hard-KD objective (Eq.~\ref{eq:loss-hard-kd}) and ignores valuable \textit{inter-class} knowledge.
    Contrarily, with privilege access to the white-box $S_0$, we can match the student's logits response $S^{\mathcal{Z}}\left(x\right)$ with that of the preceding student $S_{0}^{\mathcal{Z}}\left(x\right)$ as soft labels.
    Formally, we distill $S$ by the following composite KD objective:
    \begin{equation} \label{eq:loss-kd-total}
        \begin{aligned}
            \mathcal{L}_{S} & =\mathcal{L}_{\text{Hard-KD}}+\mathcal{L}_{\text{Soft-KD}} \\
                & = \mathbb{E}_{x\sim\mathcal{D}}\Big[\mathcal{L}_{\text{CE}}\left(\hat{y}_{S-\text{Soft}},\hat{y}_{T-\text{Hard}}\right) + \mathcal{L}_{\text{L1}}\left(S^{\mathcal{Z}}\left(x\right),S_{0}^{\mathcal{Z}}\left(x\right)\right)\Big].
        \end{aligned}
    \end{equation}
    where $\mathcal{L}_{\text{L1}}$ is the L1 divergence loss.
\end{foldable}


We summarize our method as pseudocode in Alg.~\ref{alg:dipkd}.
\begin{algorithm} % alg:dipkd
    \SetKwInput{KwParam}{Parameters}
    \SetKwComment{Comment}{}{}
    
    \caption{Proposed method \textbf{\methodname}}
    \label{alg:dipkd}

    \KwIn{a pre-trained teacher network $T$}
    \KwParam{number of synthetic images $N$}
    \KwOut{a student network $S$}

    % --------------------------------------------------------------------------
    \vspace{-0.5\baselineskip}
    \Comment{\noindent\hrulefill}

    \Comment{$\triangleright$ \textbf{\textit{Synthesis}}}
    Construct hierarchical noise images $\left\{ x_\text{hn} \right\}$ \hfill $|$ Eq.~\ref{eq:synthesis-hiernoise-multiscale},~\ref{eq:synthesis-hiernoise-combine} \\
    Apply nonlinear transforms: $\left\{ x_\text{nt}\right\} \to \left\{x_\text{hn} \right\}$ \hfill $|$ Eq.~\ref{eq:synthesis-nonlintrans} \\
    Apply semantic cutmixing: $\left\{x_\text{hn} \right\} \to \mathcal{X}_\mathcal{P}$ \hfill $|$ Eq.~\ref{eq:synthesis-diverging-filter},~\ref{eq:synthesis-semcutmix}\\
    
    \vspace{0.5\baselineskip}
    \Comment{$\triangleright$ \textbf{\textit{Contrastion}}}
    Construct $\mathcal{D} = \left\{ \hat{x} | \hat{x} \sim \mathcal{X}_\mathcal{P} \right\}_{i=1}^N$\\
    Train student $S_0$ with $T$ on $\mathcal{D}$ via $\mathcal{L}_\text{Hard-KD}$ \hfill $|$ Eq.~\ref{eq:loss-hard-kd} \\
    Initialize instance-discriminator $R$ \\
    \While{($\mathcal{D}$ not converged)} {
        Retrieve images $x\sim\mathcal{D}$ and their positive $x^+$ and negative views $x^- \ne x$ \\
        Update $x$ and $R$ via contrastive loss $\mathcal{L}_\text{CT}$ \hfill $|$ Eq.~\ref{eq:loss-contrastive}
    }
    \vspace{0.5\baselineskip}
    \Comment{$\triangleright$ \textbf{\textit{Distillation}}}
    Train student $S$ with $T$ and $S_0$ on $\mathcal{D}$ via $\mathcal{L}_{S}$ \hfill $|$ Eq.~\ref{eq:loss-kd-total}

    \vspace{0.5\baselineskip}
    \Return{$S$}
\end{algorithm}
